\documentclass[10pt, a4paper]{article}

% ── Codificação e idioma ──
\usepackage[utf8]{inputenc}
\usepackage[T1]{fontenc}
\usepackage[brazil]{babel}

% ── Layout ──
\usepackage[margin=1.4cm, top=1.2cm, bottom=1.1cm]{geometry}
\usepackage{enumitem}
\usepackage{multicol}

% ── Espaçamento ──
\setlength{\parskip}{1.5pt plus 1pt}
\setlength{\parindent}{0pt}
\setlength{\multicolsep}{2pt plus 1pt minus 1pt}

% ── Matemática ──
\usepackage{amsmath, amssymb}

% ── Tabelas e gráficos ──
\usepackage{booktabs}
\usepackage{array}
\usepackage{xcolor}
\usepackage{colortbl}
\usepackage{graphicx}
\usepackage[breakable]{tcolorbox}

% ── Cabeçalho e rodapé ──
\usepackage{fancyhdr}
\pagestyle{fancy}
\fancyhf{}
\lhead{\small Fundamentos de Sistemas Computacionais}
\rhead{\small Mapeamento de Cache}
\cfoot{\thepage}
\renewcommand{\headrulewidth}{0.4pt}

% ── Configurações de listas ──
\setlist[enumerate,1]{label=\textbf{\arabic*.}, leftmargin=*, itemsep=6pt, topsep=3pt, parsep=0pt}
\setlist[enumerate,2]{label=\alph*), leftmargin=1.5em, itemsep=2pt, topsep=3pt, parsep=0pt}
\setlist[itemize]{itemsep=1pt, topsep=2pt, parsep=0pt}

% ── Caixas de destaque ──
\newtcolorbox{dica}{
  colback=blue!5!white,
  colframe=blue!50!black,
  title={\textbf{Dica}},
  fonttitle=\small,
  sharp corners,
  boxrule=0.5pt,
  left=4pt, right=4pt, top=2pt, bottom=2pt,
  before skip=4pt, after skip=4pt
}

\newtcolorbox{formulario}{
  colback=gray!5!white,
  colframe=gray!60!black,
  title={\textbf{Formul\'{a}rio}},
  fonttitle=\small,
  sharp corners,
  boxrule=0.5pt,
  left=4pt, right=4pt, top=2pt, bottom=2pt,
  before skip=4pt, after skip=4pt
}

% ══════════════════════════════════════════
\begin{document}

\begin{center}
  {\Large \textbf{Lista de Exerc\'{i}cios}}\\[2pt]
  {\large Mapeamento Associativo e Conjunto Associativo}\\[3pt]
  {\normalsize 98800-04 -- Fundamentos de Sistemas Computacionais\ \ $\cdot$\ \ Prof.\ Ia\c{c}an\~{a} Ianiski Weber}
\end{center}

\vspace{1pt}
\hrule
\vspace{4pt}

\begin{formulario}
\small
$\text{AMAT} = \text{Hit time} + \text{Miss rate}\times\text{Miss penalty}$
\end{formulario}

\begin{enumerate}

% ── 1: Conceitual (fácil) ──
\item Em uma frase cada, explique a diferen\c{c}a entre \textbf{mapeamento direto}, \textbf{totalmente associativo} e \textbf{conjunto associativo $N$-way} quanto a \emph{em quais posi\c{c}\~{o}es da cache} um bloco da mem\'{o}ria principal pode ser alojado. Em seguida, justifique a afirma\c{c}\~{a}o: \emph{``mapeamento direto e totalmente associativo s\~{a}o apenas casos especiais do conjunto associativo''}, indicando qual valor de $N$ corresponde a cada um.

% ── 2: Bits no totalmente associativo (médio) ──
\item Considere um espa\c{c}o de endere\c{c}amento de \textbf{256\,MB}, uma cache \textbf{totalmente associativa} de \textbf{1024 linhas} e \textbf{blocos de 16 bytes}.
\begin{enumerate}
  \item Mostre a divis\~{a}o dos bits do endere\c{c}o (\emph{tag} e \emph{offset}) e explique \textbf{por que n\~{a}o existe campo de linha/\'{i}ndice} neste mapeamento.
  \item Quantos bytes de \textbf{dados \'{u}teis} a cache armazena efetivamente?
  \item Qual o tamanho (em bits) da \textbf{mem\'{o}ria associativa}, sabendo que cada entrada guarda a \emph{tag} mais um \emph{bit de validade}?
\end{enumerate}

% ── 3: Bits no conjunto associativo (médio-difícil) ──
\item Uma cache \textbf{conjunto associativa 4-way} possui \textbf{32\,KB}, \textbf{blocos de 64 bytes} e endere\c{c}os de \textbf{32 bits}.
\begin{enumerate}
  \item Quantos \textbf{blocos} cabem na cache? Quantos \textbf{conjuntos} ela possui?
  \item Determine o n\'{u}mero de bits de \textbf{offset}, \textbf{\'{i}ndice} e \textbf{tag}.
  \item Quantos \textbf{comparadores de tag} atuam \emph{simultaneamente} em uma busca? Por que esse n\'{u}mero independe do tamanho total da cache?
\end{enumerate}

% ── 4: Custo / comparadores (médio) ──
\item Para uma cache de \textbf{64\,KB} com \textbf{blocos de 4 bytes}, calcule quantos \textbf{comparadores de tag} seriam necess\'{a}rios se ela fosse \textbf{totalmente associativa} e quantos seriam necess\'{a}rios se fosse \textbf{8-way}. Use os valores para explicar por que o mapeamento conjunto associativo \emph{viabiliza caches maiores} sem estourar o custo da mem\'{o}ria associativa.

% ── 5: AMAT com L2 (difícil) ──
\item Um sistema tem \textbf{L1} (hit time $=2$ ciclos, miss rate $=8\%$) e \textbf{L2} (hit time $=12$ ciclos, miss rate $=20\%$); uma falha na L2 vai \`{a} \textbf{DRAM} com penalidade de \textbf{150 ciclos}.
\begin{enumerate}
  \item Calcule a \emph{miss penalty} da L1 e o \textbf{AMAT} do sistema (em ciclos).
  \item Recalcule o AMAT \textbf{removendo a L2} (a L1 passa a falhar direto para a DRAM). Quantas vezes a presen\c{c}a da L2 deixou o acesso m\'{e}dio mais r\'{a}pido?
\end{enumerate}

\end{enumerate}

\end{document}
